Deployed system prompts often run to thousands of tokens and dozens of rules, and prompt authors usually assume that every rule they write will be followed.
We test this assumption with a controlled, automatically scored benchmark.
It contains 40 mutually compatible rules and 4 contradictory rule pairs, and it varies system-prompt length, rule count, rule position, and rule conflicts independently under a paired design.
We run 9 open instruction-tuned models (0.5B--14.7B parameters; Qwen2.5, Llama-3.x, Qwen3), which gives 28,500 generations and 256,500 scored rule instances.
Compliance degrades continuously, and in every case more steeply than a common practitioner hypothesis predicts.
With the rule count fixed at five, per-rule compliance falls by 7.7 percentage points (pp) per doubling of prompt length beyond 500 tokens (95\% CI 7.1--8.3), about 15\% relative per doubling.
Per-rule compliance declines from five rules onward in every model, with no sharp saturation point.
Adding a contradicting rule cuts compliance on each member of a pair by 30--83\%, and for length conflicts 58--90\% of responses satisfy neither rule.
Position effects depend on structure: short rule lists show mild primacy, but in long prompts distance from the end dominates.
Moving the same five rules from the middle to the end of a 4k-token prompt raises compliance by 21\pp.
Model recipe matters more than scale: \qwenthree keeps 80\% per-rule compliance at 40 rules, while \qfourteen does no better than \qseven.
Our results suggest one concrete mitigation: put critical rules in a single block at the end of long system prompts.
